\item Las dimensiones de un salón de clase son 4.20 m $\times$ 3.00 m $\times$ 2.50 m. (a) Encuentre el número de moléculas de aire en él a presión atmosférica y $20.0^\circ\text{C}$. (b) Halle la masa de este aire, si supone que el aire consiste en moléculas diatómicas con masa molar de $28.9\text{ g/mol}$. (c) Encuentre la energía cinética promedio de las moléculas. (d) Obtenga la rapidez molecular rms. (e) \textbf{¿Qué pasaría si?} Suponga que la masa de una molécula de aire es el doble de su valor real, pero el salón contiene el mismo número de moléculas. Determine el cambio en energía interna del aire en el salón a medida que la temperatura se eleva a $25.0^\circ\text{C}$. (f) Explique cómo convencer a un estudiante de que su respuesta al inciso (e) es correcta, aun cuando sea sorprendente.
